\section{From diameter covers to projective maps}
\label{sec:geometry}

The geometric part of the proof converts a smaller-diameter cover into a
smooth map whose positive coordinates separate orthogonal lines.  We first
make this conversion explicit, so that the remaining obstruction can be
formulated entirely on projective space.

The projection-matrix metric is the rank-one case of the standard
Euclidean embedding of real Grassmannians
\cite[Section~5, Theorem~2]{ConwayHardinSloane1996}. Their chordal
distance is $1/\sqrt2$ times the Frobenius distance used here. We give
the short computation to fix the normalization throughout the proof.

Let $V=\mathbb R^k$ have its usual inner product, with $k\geq2$, and let
$\mathbb P(V)$ denote its space of one-dimensional linear subspaces.
For a line $x=[u]$, represented by a unit vector $u$, write $P_x=uu^{\mathsf T}$
for the orthogonal projector onto $x$.  Let $\operatorname{Sym}(V)$ be the
space of real symmetric operators, equipped with the Frobenius inner product
$\langle A,B\rangle_F=\operatorname{tr}(AB)$, and put
\[
 N_k=\dim\operatorname{Sym}(V)=\frac{k(k+1)}2,
 \qquad X_k=\{P_x:x\in\mathbb P(V)\}.
\]

\begin{lemma}[Projector geometry]
\label{lem:projector-geometry}
The map $x\mapsto P_x$ identifies $\mathbb P(V)$ homeomorphically with the
compact set $X_k$ in the trace-one affine subspace of
$\operatorname{Sym}(V)$.  This affine subspace is Euclidean of dimension
$N_k-1$.  If $u,v$ are unit representatives of $x,y$, then
\begin{equation}
 \|P_x-P_y\|_F^2=2-2\langle u,v\rangle^2.
 \label{eq:projector-distance}
\end{equation}
Consequently $\operatorname{diam}(X_k)=\sqrt2$, and two projectors are at
this distance exactly when their lines are orthogonal.
\end{lemma}

\begin{proof}
The formula $uu^{\mathsf T}$ is unchanged when $u$ is replaced by $-u$,
so it gives a continuous map on projective space.  Its range as an operator
is exactly $x$, which proves injectivity.  A continuous injection from the
compact space $\mathbb P(V)$ into a Hausdorff space is a homeomorphism onto
its image.  Every $P_x$ has trace one, and the trace-one hyperplane has
codimension one in $\operatorname{Sym}(V)$.  Finally,
$\operatorname{tr}(P_x^2)=\operatorname{tr}(P_y^2)=1$ and
$\operatorname{tr}(P_xP_y)=\langle u,v\rangle^2$, proving
\eqref{eq:projector-distance}.  Orthogonal lines exist because $k\geq2$.
\end{proof}

\begin{remark}[Euclidean coordinates]
An explicit isometry from the trace-one symmetric $4\times4$
matrices to $\mathbb R^9$ is
\begin{equation}
\begin{split}
 A\longmapsto\bigg(&\frac{a_{11}-a_{22}}{\sqrt2},\;
 \frac{a_{11}+a_{22}-2a_{33}}{\sqrt6},\;
 \frac{a_{11}+a_{22}+a_{33}-3a_{44}}{\sqrt{12}},\\
 &\sqrt2a_{12},\sqrt2a_{13},\sqrt2a_{14},
   \sqrt2a_{23},\sqrt2a_{24},\sqrt2a_{34}\bigg).
\end{split}
\label{eq:euclidean-coordinates}
\end{equation}
The first three coordinates use an orthonormal basis of the trace-zero
diagonal matrices.  The factors $\sqrt2$ in the remaining six coordinates
account for the two equal off-diagonal entries in the Frobenius norm.
Thus the nine-dimensional witness will be $X_4$ with exactly its stated
metric.
\end{remark}

Write $[m]=\{1,\ldots,m\}$ and
$\Delta^{m-1}=\{(t_1,\ldots,t_m):t_i\ge0,\ \sum_i t_i=1\}$.

\begin{definition}
\label{def:admissible}
A smooth map $f=(f_1,\ldots,f_m):\mathbb P(V)\to\mathbb R^m$ is
\emph{admissible} if each $f_i$ is nonnegative, $\sum_i f_i=1$, and,
writing $S(x)=\{i:f_i(x)>0\}$,
\begin{equation}
 x\perp y\quad\Longrightarrow\quad S(x)\cap S(y)=\varnothing.
 \label{eq:admissibility}
\end{equation}
Here smoothness means smoothness of the real coordinate functions on the
usual smooth manifold $\mathbb P(V)$.
\end{definition}

Restricting an admissible map to $\mathbb P(W)$ for a nonzero linear
subspace $W\subseteq V$ preserves admissibility.  Coordinates that vanish
identically on this subspace may be omitted.

\begin{lemma}[The strict-gap reduction]
\label{lem:cover-to-map}
If $X_k$ is covered by at most $m$ subsets, each of diameter strictly less
than $\sqrt2$, then there is an admissible map
$f:\mathbb P(V)\to\mathbb R^m$.
\end{lemma}

\begin{proof}
Pad the cover by empty sets to write it as $A_1,\ldots,A_m$.
For every nonempty $A_i$, set $d_i=\operatorname{diam}(A_i)$ and
$\varepsilon_i=(\sqrt2-d_i)/4>0$, and define the relatively open set
\[
 U_i=\{P\in X_k:\operatorname{dist}_F(P,A_i)<\varepsilon_i\}.
\]
Set $U_i=\varnothing$ if $A_i=\varnothing$.  The sets $U_i$ cover $X_k$.
For any $P,Q\in U_i$, points of $A_i$ within $\varepsilon_i$ of $P,Q$
and the triangle inequality give
\[
 \|P-Q\|_F<d_i+2\varepsilon_i
 =\frac{d_i+\sqrt2}{2}<\sqrt2.
\]
In particular $\operatorname{diam}(U_i)\leq(d_i+\sqrt2)/2<\sqrt2$.

Pull these sets back to a finite open cover $V_i$ of $\mathbb P(V)$.
A smooth partition of unity subordinate to this cover
\cite[Theorem~2.23]{Lee2013} gives smooth
nonnegative functions $f_i$, with sum one and $f_i(x)>0$ only if
$x\in V_i$.  Empty sets receive the zero function.  If $i$ belonged to
both $S(x)$ and $S(y)$ for orthogonal lines $x,y$, then $P_x,P_y\in U_i$
would be at distance $\sqrt2$ by Lemma~\ref{lem:projector-geometry},
a contradiction.
\end{proof}

The strict diameter inequalities are what permit the open neighborhoods
and the smooth partition of unity; merely excluding individual orthogonal
pairs from arbitrary color classes does not supply this reduction.

For an admissible map, define its \emph{support complex} $K$ on the label
set $\{1,\ldots,m\}$ by
\[
 I\in K\quad\Longleftrightarrow\quad
 I\subseteq S(x)\ \text{for some }x\in\mathbb P(V).
\]
It is a finite simplicial complex, possibly with unused labels, and $f$
takes values in its geometric realization inside the standard simplex.
No triangulation of $\mathbb P(V)$ is assumed.  The next sections derive
the properties of $K$ that will rule out an admissible map when $k=4$ and
$m=10$.
